\section{El sistema completo: población, sustrato y ambiente controlado}

\subsection{De una larva a una población}

El modelo de Eriksen describe la dinámica de los compartimentos de una larva
individual. Un sistema de bioconversión real contiene una población de
$N(t)$ larvas, cuyas tasas agregadas determinan los flujos de materia y
energía a escala del contenedor.

\begin{definition}[Tasa agregada de población]\label{def31}
Sea $r(t)$ una tasa por larva (asimilación, síntesis, producción de CO$_2$)
y $N(t)$ el número de larvas vivas en el contenedor. La \emph{tasa agregada}
de la población es
\begin{equation}
R(t) = N(t)\,r(t),
\label{eq:agregada}
\end{equation}
bajo el supuesto de población homogénea: todas las larvas comparten el mismo
estado $(A, B, L)$, es decir, la misma edad, peso y condición fisiológica.
\end{definition}

\begin{remark}\label{rem:homogeneidad}
El supuesto de población homogénea es razonable en crías sincronizadas
(larvas neonatas sembradas el mismo día, práctica estándar en producción),
pero se degrada conforme avanza el ciclo: la variabilidad individual en las
tasas de crecimiento genera una distribución de pesos que se ensancha con el
tiempo. Un refinamiento posible ---que dejamos como extensión--- es
reemplazar \eqref{eq:agregada} por la integral sobre la distribución de
pesos de la población, $R(t) = \int r(w)\,n(w,t)\,dw$, a costa de convertir
el modelo en una ecuación de estructura poblacional. En lo que sigue se
adopta \eqref{eq:agregada}, y la mortalidad se admite mediante $N(t)$
decreciente si los datos lo requieren.
\end{remark}

\subsection{El contenedor como volumen de control}

\begin{definition}[Sistema de bioconversión instrumentado]\label{def32}
El \emph{sistema completo} es el volumen de control delimitado por las
paredes del contenedor de cría, cuya frontera es atravesada por dos flujos
de aire ---\emph{inlet} y \emph{outlet}--- y en $t = 0$ por la carga inicial
del lote (larvas neonatas y sustrato). Su vector de estado se organiza en
tres grupos:
\begin{equation}
\mathbf{x} =
\underbrace{(A,\, B,\, L)}_{\text{larva promedio}}
\;\cup\;
\underbrace{(N,\, S)}_{\text{población y sustrato}}
\;\cup\;
\underbrace{(T,\, w,\, c_{\mathrm{CO_2}},\, c_{\mathrm{CH_4}},\, c_{\mathrm{O_2}})}_{\text{ambiente del aire}},
\label{eq:estado-completo}
\end{equation}
donde $S$ es la masa de sustrato disponible, $T$ la temperatura del aire,
$w$ la humedad absoluta y $c_i$ las concentraciones de los gases medidos.
\end{definition}

\begin{remark}\label{rem:frontera}
La elección de la frontera determina qué es flujo y qué es transformación
(cf.~Fig.~\ref{fig:balance}). Con la frontera en las paredes del contenedor,
la ventilación es el único flujo continuo que cruza la frontera
($J_{\mathrm{in}}$, $J_{\mathrm{out}}$), mientras que el consumo de
sustrato, el crecimiento larval, la respiración, la evaporación y la
metanogénesis son transformaciones internas $P$, $C$. El aire del
\emph{inlet} aporta O$_2$, vapor de agua y calor sensible a condiciones
controladas $(T_{\mathrm{in}}, w_{\mathrm{in}})$; el \emph{outlet} remueve
aire empobrecido en O$_2$ y enriquecido en CO$_2$, CH$_4$ y vapor, y es
precisamente en esa corriente de salida donde se instalan las mediciones
$\mathbf{y} = (T, H, \mathrm{CO_2}, \mathrm{CH_4}, \mathrm{O_2})$.
\end{remark}

\begin{remark}[Operación por lote]\label{rem:batch}
El sistema opera \emph{por lote}: en $t = 0$ se siembra una población de
$N_0$ larvas neonatas sobre una masa inicial $S_0$ de sustrato, se cierra el
contenedor y el ciclo transcurre sin recargas hasta que las larvas alcanzan
el estado de prepupa. No hay eventos discretos de \emph{carga}: el estado
$\mathbf{x}(t)$ evoluciona continuamente en todo el horizonte $[0, t_f]$.
La única conmutación programada es la de la ventilación: cada dos días el
\emph{inlet} se cierra durante un intervalo $[t_k^{c},\, t_k^{c}+\Delta]$
para estimar las tasas de producción de gases por acumulación (método de
cámara estática, Sección~6). Estas conmutaciones no alteran el estado, solo
activan o desactivan los términos de flujo de los balances
(cf.~Sección~4.2). Además de las mediciones de gas, el contenedor
está montado sobre una celda de carga que registra la masa total del lote,
$m_{\mathrm{LC}}(t) = S(t) + N(t)\bigl(A(t) + B(t) + L(t)\bigr) + m_f$,
donde $m_f$ es la tara (contenedor y estructura). La tasa de pérdida de masa
del lote obedece al flujo neto a través de la frontera:
\begin{equation}
-\dot{m}_{\mathrm{LC}} \;=\;
\dot{m}_{\mathrm{CO_2}} + \dot{m}_{\mathrm{CH_4}} + \dot{m}_{\mathrm{H_2O}}
\;-\; \dot{m}_{\mathrm{O_2,in}},
\label{eq:celda-carga}
\end{equation}
es decir, la masa de CO$_2$, CH$_4$ y vapor de agua que sale por el
\emph{outlet}, menos la masa de O$_2$ (y demás componentes) que el aire del
\emph{inlet} aporta y queda incorporada en esos productos. La
Ec.~\eqref{eq:celda-carga} es una restricción de consistencia global que
servirá de validación independiente de los balances de gases.
\end{remark}

\begin{remark}[Actividad microbiana del sustrato]\label{rem:microbiano}
El sustrato no es un almacén pasivo de alimento: alberga una microbiota que
respira ---consume O$_2$ y excreta CO$_2$---, genera calor y, en zonas
anaerobias del lecho húmedo, produce CH$_4$ por metanogénesis. En
consecuencia, las señales medidas en el \emph{outlet} son superposiciones de
la actividad larval y la microbiana,
\begin{equation}
r_{\mathrm{CO_2}}^{\mathrm{medido}} = N\,r_{\mathrm{CO_2}}
+ r_{\mathrm{CO_2}}^{\mathrm{mic}}(S, T, w),
\qquad
r_{\mathrm{O_2}}^{\mathrm{medido}} = N\,r_{\mathrm{O_2}}
+ r_{\mathrm{O_2}}^{\mathrm{mic}}(S, T, w),
\label{eq:co2-superpuesto}
\end{equation}
donde $r_{\mathrm{O_2}}^{\mathrm{medido}}$ denota la tasa total de consumo
de oxígeno (cf.~Ec.~\eqref{eq:rq}). El CH$_4$ proviene casi exclusivamente
del segundo proceso, lo que lo convierte en un indicador indirecto de la
anaerobiosis del sustrato. Nótese, además, que el cociente respiratorio
medido en el \emph{outlet} no es el cociente larval puro, sino
\begin{equation}
RQ_{\mathrm{medido}} \;=\;
\frac{N\,r_{\mathrm{CO_2}} + r_{\mathrm{CO_2}}^{\mathrm{mic}}}
     {N\,r_{\mathrm{O_2}} + r_{\mathrm{O_2}}^{\mathrm{mic}}},
\label{eq:rq-medido}
\end{equation}
de modo que la separación de ambas contribuciones ---un problema de
identificabilidad que se abordará en la sección de mediciones--- es
necesaria para interpretar $RQ$ en términos fisiológicos. Una vía
experimental directa es un lote testigo sin larvas, que mide
$r_{\mathrm{CO_2}}^{\mathrm{mic}}$ y $r_{\mathrm{O_2}}^{\mathrm{mic}}$ de
forma independiente.
\end{remark}

\begin{figure}[htbp]
\centering
\begin{tikzpicture}[
  font=\small,
  cont/.style={draw, very thick, rounded corners=6pt, minimum width=11.0cm,
               minimum height=4.8cm, fill=gray!5},
  flujo/.style={-{Stealth[length=3mm]}, very thick},
  comp/.style={draw, thick, rounded corners=3pt, fill=blue!8,
               minimum width=3.3cm, minimum height=1.55cm, align=center},
  sens/.style={draw, thick, fill=red!12, rounded corners=2pt,
               minimum width=1.0cm, minimum height=0.65cm, align=center},
  lbl/.style={fill=white, inner sep=1.6pt, font=\footnotesize, align=center},
  lblf/.style={above=2pt, midway, font=\footnotesize, align=center}
]

% Contenedor
\node[cont] (box) {};
\node[anchor=north west, xshift=8pt, yshift=-6pt] at (box.north west)
  {\textbf{Contenedor (volumen de control)}};

% Compartimentos internos (más arriba)
\node[comp] (sus) at ([xshift=-3.6cm,yshift=-0.4cm]box.center)
  {Sustrato\\[2pt] $S$};
\node[comp, fill=orange!12] (lar) at ([xshift=0cm,yshift=-0.4cm]box.center)
  {Población larval\\[2pt] $N$\\ $(A,B,L)$};
\node[comp, fill=green!8] (air) at ([xshift=3.6cm,yshift=-0.4cm]box.center)
  {Aire interno\\[2pt] $T,\, w$\\ $c_{\mathrm{CO_2}},\, c_{\mathrm{CH_4}},\, c_{\mathrm{O_2}}$};

% Intercambios internos: sustrato <-> larvas
\draw[{Stealth[length=2.2mm]}-{Stealth[length=2.2mm]}, thick, orange!85!black]
  (sus) -- node[lbl, above=6pt, midway]{consumo} (lar);

% Larvas -> aire: CO2, calor, vapor (flecha superior)
\draw[-{Stealth[length=2.2mm]}, thick, orange!85!black]
  (lar.15) -- node[lbl, above=6pt, midway]{CO$_2$, calor, vapor} (air.165);
% Aire -> larvas: O2 (flecha inferior)
\draw[-{Stealth[length=2.2mm]}, thick, orange!85!black]
  (air.195) -- node[lbl, below=6pt, midway]{O$_2$} (lar.345);

% Sustrato -> aire: evaporación, CH4, CO2 microbiano (curva inferior)
\draw[{Stealth[length=2.2mm]}-{Stealth[length=2.2mm]}, thick, orange!85!black, dashed]
  (sus.south) .. controls ([xshift=1.4cm,yshift=-1.1cm]sus.south)
    and ([xshift=-2.0cm,yshift=-1.2cm]air.south) ..
  node[lbl, pos=0.5]{evaporación, CH$_4$, CO$_2$ mic. / O$_2$ mic.} (air.south);

% Inlet (rótulo a la izquierda, dos líneas, flecha despejada)
\draw[flujo, green!50!black] ([yshift=0.9cm]box.west) ++(-2.0,0)
  -- node[pos=0, above=2pt, anchor=south west, font=\footnotesize, align=left]
     {inlet: $T_{\mathrm{in}},\; w_{\mathrm{in}}$,\\ $c_{\mathrm{O_2,in}},\; Q$}
  ([yshift=0.9cm]box.west);

% Outlet (rótulo a la derecha, flecha despejada)
\draw[flujo, red!70!black] ([yshift=1.3cm]box.east)
  -- node[pos=1, above=2pt, anchor=south east, font=\footnotesize, align=right]
     {outlet: $Q$}
  ++(2.0,0);

% Sensores en el outlet (debajo de la flecha)
\node[sens] (sensores) at ([xshift=1.0cm,yshift=0.55cm]box.east) {$\mathbf{y}$};
\draw[-{Stealth[length=2mm]}, thick, red!70!black, dashed]
  (sensores.north) -- ([xshift=1.0cm]box.east);
\node[below=3pt of sensores, align=center, font=\footnotesize]
  {$T,\ H,\ \mathrm{O_2}$\\ $\mathrm{CO_2},\ \mathrm{CH_4}$};

\end{tikzpicture}
\caption{El sistema completo de bioconversión instrumentado
(Definición~\ref{def32}). Dentro del volumen de control interactúan tres
grupos de compartimentos: sustrato, población larval y aire interno; las
flechas indican el sentido de cada intercambio (el O$_2$ fluye del aire
hacia larvas y microbiota; el CO$_2$, el calor, el vapor y el CH$_4$, hacia
el aire). La frontera es atravesada únicamente por los flujos de ventilación
(\emph{inlet/outlet}, caudal $Q$); las mediciones $\mathbf{y}$ se toman en
la corriente de salida y una celda de carga registra la masa total del
lote.}
\label{fig:sistema-completo}
\end{figure}

\subsection{Condiciones iniciales y horizonte del ciclo}

\begin{definition}[Ciclo de bioconversión]\label{def33}
Un \emph{ciclo de bioconversión} es la trayectoria del sistema
\eqref{eq:estado-completo} en el horizonte $[0, t_f]$, con condiciones
iniciales
\begin{equation}
N(0) = N_0, \quad S(0) = S_0, \quad
A(0) = A_0, \quad B(0) = B_0, \quad L(0) = L_0,
\label{eq:cond-iniciales}
\end{equation}
y ambiente en equilibrio con el \emph{inlet},
$T(0) = T_{\mathrm{in}}$, $w(0) = w_{\mathrm{in}}$,
$c_i(0) = c_{i,\mathrm{in}}$. El ciclo termina en $t_f$, definido como el
instante en que la fracción de prepupas en la población supera un umbral
$\pi_f$ convenido (típicamente $\pi_f \approx 0.5$), o equivalentemente
cuando la tasa agregada de asimilación cae por debajo del umbral de
mantenimiento de la población,
\begin{equation}
t_f \;=\; \inf\bigl\{ t > 0 \;:\; R_A(t) \;\leq\; N(t)\,m\,B(t) \bigr\}.
\label{eq:tf}
\end{equation}
En las condiciones de cría previstas, $t_f$ se espera dentro de la ventana
de $14$ a $16$ días. Cada ciclo se replica $n_{\mathrm{rep}} = 3$ veces en
condiciones nominales idénticas, lo que permitirá estimar la variabilidad
entre lotes de los parámetros del modelo (Sección~6).
\end{definition}

\begin{remark}[Calendario de cierres de ventilación]\label{rem:cierres}
El protocolo experimental incluye cierres del \emph{inlet} cada dos días,
en los instantes $t_k^{c} = 2k$ días, $k = 1, 2, \dots$, cada uno de
duración $\Delta$ por estimar (Sección~6). Durante un cierre, $Q = 0$ y las
concentraciones del aire interno evolucionan únicamente por la producción
interna, de modo que las pendientes $\Delta c_i / \Delta t$ miden
directamente las tasas $R_i$. El calendario $\{t_k^{c}\}$ y la duración
$\Delta$ son parte del protocolo: se conocen con exactitud y se registran
junto con las mediciones.
\end{remark}

\begin{remark}[Criterios de terminación y su medición]\label{rem:tf}
Los dos criterios de \eqref{eq:tf} no son equivalentes operativamente: la
fracción de prepupas requiere inspección visual o muestreo, mientras que el
criterio metabólico es \emph{observable en línea} ---con la salvedad de la
contribución microbiana al CO$_2$ (Observación~\ref{rem:microbiano}), que no
cesa al final del ciclo. Una firma de terminación más robusta es el cambio
de régimen del cociente CH$_4$/CO$_2$: al llegar la prepupa, las larvas
dejan de alimentarse y de remover el sustrato, cesa la bioturbación que
aireaba el lecho, se expanden las zonas anaerobias y el CH$_4$ aumenta
mientras el CO$_2$ larval cae a su componente de mantenimiento. Esta firma
es la que el gemelo digital puede evaluar en tiempo real sin intervenir el
lote. Para larvas neonatas de cría sincronizada, los valores iniciales
$A_0, B_0, L_0$ corresponden al peso seco promedio de una larva de un día,
repartido entre compartimentos según su composición bioquímica medida.
\end{remark}

\paragraph{Variables monitoreadas fuera del vector de estado.}
El mapa de observación $\mathbf{y}$ definido hasta aquí se restringe a las
variables que participan directamente en los balances del modelo:
$T$, $H$, CO$_2$, CH$_4$ y O$_2$. La instrumentación del sistema puede,
sin embargo, incluir sensores adicionales ---notablemente H$_2$S y NH$_3$---
cuyas señales no se incorporan como variables de estado. La razón es de
economía de modelado: incluir el H$_2$S exigiría abrir el balance de azufre
del sustrato, y el NH$_3$, el balance de nitrógeno con su equilibrio
amonio--amoníaco dependiente de pH; en ambos casos el costo en compartimentos
y parámetros adicionales es desproporcionado frente al valor que aportan al
objetivo central, que es describir la bioconversión de materia y energía.
Estos gases se clasifican entonces como \emph{variables de supervisión}:
mediciones en tiempo real que informan sobre la salud del proceso sin
alimentar las ecuaciones de estado.

\paragraph{Su rol en la capa de supervisión del gemelo digital.}
Aunque fuera del modelo dinámico, estas señales tienen una función bien
definida en el sistema monitoreado. El H$_2$S es un indicador de
putrefacción anaerobia de la fracción proteica del sustrato: su aparición
señala condiciones de falla ---exceso de humedad, compactación del lecho o
sustrato proteico sin consumir--- y opera como alarma temprana sobre el lote.
El NH$_3$, por su parte, es la vía gaseosa de pérdida de nitrógeno del
sistema: su concentración informa sobre la volatilización de amonio y,
por tanto, sobre la calidad final del frass como fertilizante. La distinción
introducida aquí ---variables de \emph{estado} que el modelo explica versus
variables de \emph{supervisión} que el sistema vigila--- es un principio
general del diseño: la frontera del modelo no la define lo que se puede
medir, sino lo que el modelo se propone explicar. Las extensiones a los
balances de nitrógeno y azufar quedan planteadas como trabajo futuro, una
vez validado el núcleo materia--energía del sistema.